% corpus/test/supervision_alertes.tex
\documentclass[11pt,a4paper]{report}
\usepackage[utf8]{inputenc}
\usepackage[T1]{fontenc}
\usepackage[french]{babel}

\begin{document}

\chapter{Supervision et alertes}

\emph{Document de test fictif, rédigé pour valider le RAG navigation. Cellule Exploitation (E1), version 1.0 du 22 octobre 2026.}

\section{Outils de supervision}
La supervision de GANDAL est centralisée sur la machine virtuelle VM160. Elle repose sur les outils suivants.

\begin{itemize}
\item Prometheus collecte les métriques, sur le port 9090.
\item Grafana affiche les tableaux de bord, sur le port 3000, à l'adresse \texttt{http://grafana.gandal.lan:3000}.
\item Alertmanager envoie les notifications, sur le port 9093.
\item node\_exporter tourne sur chaque machine et expose ses métriques sur le port 9100.
\item dcgm-exporter tourne sur le nœud REM et expose les métriques des cartes graphiques sur le port 9400.
\end{itemize}

Les métriques sont conservées 30 jours dans Prometheus.

\section{Tableaux de bord}
Grafana propose quatre tableaux de bord principaux : « Vue générale » pour l'état de toutes les machines, « Cluster NOVA » pour la charge des hyperviseurs, « GPU REM » pour l'utilisation, la mémoire et la température des cartes graphiques, et « Services » pour la disponibilité de GANDAL-AI, de LiteLLM et d'Ollama.

Tout membre du projet peut consulter les tableaux de bord avec son compte GANDAL. Seule la cellule Exploitation peut les modifier.

\section{Règles d'alerte}
\begin{tabular}{|l|l|l|}
\hline
Alerte & Condition & Gravité \\
\hline
Machine injoignable & aucune réponse pendant 2 minutes & critique \\
Processeur saturé & plus de 90 \% pendant 10 minutes & avertissement \\
Disque presque plein & plus de 85 \% d'occupation & avertissement \\
Disque plein & plus de 95 \% d'occupation & critique \\
GPU trop chaud & plus de 85 °C pendant 5 minutes & avertissement \\
Sauvegarde échouée & échec d'une tâche sur ARCHIVE & critique \\
GANDAL-AI dégradé & route /health en état degraded pendant 5 minutes & avertissement \\
\hline
\end{tabular}

\section{Notifications}
Les avertissements sont publiés sur le canal Mattermost \texttt{\textasciitilde{}gandal-alertes}. Les alertes critiques sont publiées sur le même canal et envoyées en plus par SMS à la personne d'astreinte. Une alerte reste active tant que sa condition est vraie ; elle est notifiée de nouveau toutes les quatre heures si personne ne l'a acquittée.

\section{Astreinte}
La cellule Exploitation assure une astreinte tournante d'une semaine. La relève a lieu le lundi à 8 h. Le nom et le numéro de la personne d'astreinte sont affichés en tête du canal \texttt{\textasciitilde{}gandal-alertes}.

En heures ouvrées, la personne d'astreinte intervient dans les 30 minutes. La nuit et le week-end, le délai d'intervention est de 2 heures.

\section{Ajouter une machine à la supervision}
\begin{enumerate}
\item Installer node\_exporter sur la machine et vérifier qu'il répond sur le port 9100.
\item Ouvrir le port 9100 depuis VM160 en faisant une demande à la cellule Réseau.
\item Ajouter la machine dans le fichier \texttt{targets.yml} de Prometheus sur VM160.
\item Vérifier dans Grafana que la machine apparaît dans le tableau de bord « Vue générale ».
\end{enumerate}

\end{document}
